% !TEX root = main.tex

% \cite{li2016unsupervised} uses a graph of nearest neighbors (NN) on SIFT descriptors \cite{lowe1999object}.

\noindent\textbf{Unsupervised learning of features.}
Several approaches related to our work learn deep models with no supervision.
Coates and Ng~\cite{coates2012learning} also use $k$-means to pre-train convnets, but learn each layer sequentially
in a bottom-up fashion, while we do it in an end-to-end fashion.
Other clustering losses~\cite{yang2016joint,xie2016unsupervised,dosovitskiy2014discriminative,liao2016learning} 
have been considered to jointly learn 
convnet features and image clusters but they have never been tested on a scale to allow a thorough study on modern convnet architectures.
Of particular interest,
Yang~\etal~\cite{yang2016joint} iteratively learn convnet features and clusters with a recurrent framework.
Their model offers promising performance on small datasets but may be 
challenging to scale to the number of images required for convnets to be competitive.
Closer to our work, Bojanowski and Joulin~\cite{bojanowski2017unsupervised} learn visual features on a large dataset with
a loss that attempts to preserve the information flowing through the network~\cite{linsker1988towards}. 
Their approach discriminates between images in a similar way as examplar SVM~\cite{malisiewicz2011ensemble}, while we are simply clustering them.
\\

\noindent\textbf{Self-supervised learning.}
A popular form of unsupervised learning, called ``self-supervised learning''~\cite{de1994learning},
uses pretext tasks to replace the labels annotated by humans by ``pseudo-labels'' directly computed from the raw input data.
For example,
Doersch~\etal~\cite{doersch2015unsupervised} use the prediction of the relative position of patches in an image as a pretext task,
while Noroozi and Favaro~\cite{noroozi2016unsupervised} train a network to spatially rearrange shuffled patches.
Another use of spatial cues is the work of Pathak~\etal~\cite{pathak2016context} where missing pixels are guessed based
on their surrounding.
Paulin~\etal~\cite{paulin2015local} learn patch level Convolutional Kernel Network~\cite{mairal2014convolutional} using an image retrieval setting. 
Others leverage the temporal signal available in videos by 
predicting the camera transformation between consecutive frames~\cite{agrawal2015learning}, exploiting the temporal coherence of tracked patches~\cite{wang2015unsupervised} or segmenting video based on motion~\cite{pathak2016learning}. 
Appart from spatial and temporal coherence, many other signals have been explored: image colorization~\cite{zhang2016colorful,larsson2016learning}, cross-channel prediction~\cite{zhang2016split}, sound~\cite{owens2016ambient} or instance counting~\cite{noroozi2017representation}.
More recently, several strategies for combining multiple cues have been proposed~\cite{wang2017transitive,doersch2017multi}.
Contrary to our work, these approaches are domain dependent, requiring expert knowledge 
to carefully design a pretext task that may lead to transferable features.
\\

\noindent\textbf{Generative models.}
Recently, unsupervised learning has been making a lot of progress on image generation.
Typically, a parametrized mapping is learned between a predefined random noise and the images,
with either an autoencoder~\cite{kingma2013auto,bengio2007greedy,huang2007unsupervised,masci2011stacked,vincent2010stacked}, 
a generative adversarial network (GAN)~\cite{goodfellow2014generative} 
or more directly with a reconstruction loss~\cite{bojanowski2017optimizing}.
Of particular interest, 
the discriminator of a GAN can produce visual features, but their performance are relatively disappointing~\cite{donahue2016adversarial}. 
Donahue~\etal~\cite{donahue2016adversarial} and Dumoulin~\etal~\cite{dumoulin2016adversarially} have shown that adding an encoder
to a GAN produces visual features that are much more competitive.

